\section{交叉验证评估的是完整训练流程}
\label{sec:11-Mathematical-Statistics/06-Resampling-and-Model-Selection/03}

假设你手上有 200 个肺癌患者的回顾性数据，50 个候选生物标志物。目标是训练一个 logistic 回归，预测术后并发症风险，用 AUC 向临床科室报告模型表现。操作顺序看起来顺理成章：先对所有 200 人统一做标准化，筛掉方差过低的变量，再按单变量 p 值挑出 15 个最强标志物，最后跑一个 10 折交叉验证——AUC 0.82，还不错。直到审稿人问了一句：``特征筛选是在每个折里面做的，还是先用了全部数据？''

你重新跑一遍：每个折内独立完成标准化和变量筛选，再训练模型。AUC 变成 0.76。0.06 的差距不是一个随机波动——它是验证折信息通过全体数据的变量排序，提前进入了模型选择。

这就是交叉验证最容易被误解的地方。公式本身很简单，但把它用对的代价是理解它评估的到底是什么。

\(K\) 折交叉验证把 \(n\) 个样本分成 \(K\) 个尽可能等大的互斥折。对每个折 \(k=1,\ldots,K\)，用其余 \(K-1\) 折训练模型，在第 \(k\) 折上计算损失。最后把每个样本在它被排除的那一轮中的损失平均：

\[
\widehat R_{\mathrm{CV}} = \frac1n
\sum_{i=1}^n
L
\bigl(
y_i,
\hat f^{(-k(i))}(x_i)
\bigr),
\]

其中 \(\hat f^{(-k(i))}\) 表示没有使用第 \(i\) 个样本所在折训练的模型。每个样本恰好当一次验证，每个折恰好被排除一次。

\subsection{被评估的是训练流水线，不是一个参数向量}

很多人在写完公式后会自然地想象：CV 在评估那个最终在全数据上训练出的模型。这个直觉刚好反了。

CV 的每一折都从零开始走完整流程——预处理、特征工程、模型拟合，全部只碰训练折的数据。因此 CV 分数估计的是：``如果你只有约 \(n(K-1)/K\) 个样本，按照这个协议走完全部步骤，得到的预测器在同样分布的新数据上大概表现如何''。它不是对最终那个包含了全部 \(n\) 个样本的参数向量的无偏估计——训练样本量差了约 \(n/K\) 个，这一差异会带来系统偏差。

更进一步看，不同折的训练集高度重叠（任意两折共享 \(n(K-2)/K\) 个样本），导致各折损失之间存在强相关。把 \(K\) 个折分数当作 \(K\) 个独立观测去算标准误，通常会显著低估不确定性。\(K\) 取大值让训练规模接近全数据，偏差减小，但计算代价上升且折间相关性更强；\(K\) 取小值让验证折更大，方差估计更稳定，但训练样本减少带来的偏差更明显。5 折和 10 折是实践中跑得动的折中，不是定理要求的最优解，也不存在对所有问题的统一最优 \(K\)。

\subsection{预处理的信息泄漏比看上去隐蔽}

先标准化再 CV，绝大多数人都知道是错的。但更隐蔽的泄漏长这样：

你先把所有 200 个样本的 50 个变量各自计算均值和中位数，发现某几个变量缺失率超过 40\%，于是直接删除；又发现某变量的 99 分位值偏离均值 8 个标准差，于是按此阈值截尾。做完这些``清洗''后，你再把数据送进 CV。问题是：缺失率和异常值阈值的判断依据来自全部 200 人——你已经用验证折的信息决定了哪些变量存活、哪些值被视为异常。

正确的流程是在每个折内部独立完成整套预处理，不得越过折边界查看任何统计量：

\begin{enumerate}
\tightlist
\item
  仅用训练折的数据拟合所有预处理参数（均值、标准差、中位数填充值、PCA 方向、特征选择阈值）；
\item
  用这些参数变换训练折，也用同样的参数变换验证折——验证折不能参与任何参数估计；
\item
  在变换后的训练折上拟合模型；
\item
  在变换后的验证折上计算损失。
\end{enumerate}

这里的关键判断标准不是``有没有直接把验证标签送进损失函数''。特征选择时你并没有看验证标签，但你看了验证样本的自变量分布，通过全体数据的变量排序决定了哪些变量进入模型——验证折的信息已经渗透进了模型结构。类不平衡时的 SMOTE 过采样、早停的 epoch 选择、异常值检测的阈值，全部属于训练流水线的一部分，全部必须在折内完成。

\subsection{调参和评估用同一组 CV 分数会乐观}

如果你跑 10 折 CV 去比较 5 个不同的正则化强度 \(\lambda\)，然后挑出平均 AUC 最高的那个 \(\lambda^*\)，把该轮的 0.82 报告为泛化性能——你又重复了同一个错误。选择过程会利用验证折上的噪声：总有某个 \(\lambda\) 因为随机波动恰好在这组折上表现好，不代表它在真正的新数据上也好。

Nested cross-validation 把调参和评估的折分开：

\begin{itemize}
\tightlist
\item
  外层折定义哪些样本从未参与任何选择；
\item
  内层折在外层的训练部分中运行完整的调参流程（网格搜索、Bayesian optimization 等），选出最优配置；
\item
  外层在未见折上评估``选了最优配置 + 重新训练''这一完整协议。
\end{itemize}

外层分数是对``包含调参在内的整个建模流程''的泛化估计。内层选出的最优超参数可能随外层折不同而变化——如果某个变量在五折中只有三折被选中，你就知道它的信号强度处在可被噪声淹没的边界。

最终你还可以在全数据上按同样的调参规则选一次超参数并重训，但那个最终模型的泛化性能只能靠之前嵌套 CV 的外层分数来估计，或者靠一个从未被查看过的独立测试集。测试集一旦被用来修改模型，不管是改超参数、换特征还是调整预处理，它就退化为验证集，不再给出无偏估计。一个项目只允许一次真正的测试集评估。

\subsection{切分方式必须反映预测的实际场景}

随机打乱后按行切分，隐含了一个强假设：任意两个样本可以交换位置而不改变数据生成机制。时间序列显然不满足：用 2019 年的数据去``预测''2018 年，等于让未来泄漏到过去。还有一些同样容易被忽略的依赖结构：

\begin{itemize}
\tightlist
\item
  同一个患者的多条随访记录必须分到同一折，否则模型可以通过记忆患者身份来降低验证误差，上线后面对陌生患者时失效。
\item
  来自同一家医院、同一台 CT 设备或多中心试验中同一研究中心的样本，构成一个自然簇。随机逐行切分会把同一簇的样本散布在训练和验证之间，评估出的泛化能力严重高估——更换医院后模型可能直接崩溃。
\item
  稀有类别可以通过分层抽样保持训练和验证中的比例，但分层不能跨越已经存在的簇结构。先按患者分组，再在组内分层维持类别比例，两个约束同时生效。
\item
  训练数据来自 A 城市的三甲医院，部署目标是 B 县城的社区卫生中心。此时随机 CV 给出的低误差没有保证。验证集的分布应尽可能模拟目标部署域，而不是训练集的混合分布。
\end{itemize}

交叉验证只能估计与切分机制相符的泛化风险。如果切分时的可交换性假设在未来部署中消失——人群特征变了、设备换了、时间推移了——CV 分数就只是过去数据自洽性的度量，不是未来表现的预测。

\subsection{什么使 CV 分数不可靠}

除了信息泄漏和切分方式，还有几个因素会破坏 CV 分数的可信度：

折间的依赖结构。10 折 CV 产生 10 个分数，但它们背后只有一组数据。把这 10 个数当独立样本算标准误给出的置信区间过窄。更诚实的做法是报告逐样本 out-of-fold 损失分布、多次随机重复 CV 的分布，或比较两个算法时在同一组折上做配对差异检验。

指标的选择本身也构成一个选择步骤。平方误差、绝对误差、log loss、AUC 和校准误差衡量完全不同的目标。如果你的分析流程是``先跑 5 个指标，选最好看的那个写进文章''，那指标选择又引入了一重乐观偏差。训练时的损失函数、调参时的评价指标和最终汇报的业务指标可以不同，但它们之间的关系——以及你为什么相信优化前两个会改善第三个——应该在分析计划中预先说明。一个常见陷阱：不平衡分类问题中准确率可以高达 95\% 却完全没用，因为模型只需全部预测多数类；概率模型只看 AUC 而忽略校准误差，可能排出的序还行，但输出的``概率''偏差十几二十个百分点。

CV 分数的有限性本身也值得重视。当数据集只有 50 个样本时，5 折 CV 的每折训练集只有 40 个样本。重抽样得到的分数分布在这么小的数据量下可能非常宽——0.76 旁边可能是 0.68 和 0.83。报告``CV 均值 0.82''时的小数点后几位不是精度，是假象。模型选择的不可消除的不确定性和数据集的固有有限性，不能被一个多位数均值和一个过窄的误差棒掩盖。
